% =====================================================================
%  Professioni 6 — Il 3D artist
%  Fonte: trascrizione della presentazione del corso di 3D art ai Dev
%  Days (AIV6): un incontro online su Zoom di circa un'ora, tenuto da
%  Giacomo Caneschi, coordinatore del corso di 3D art di AIV. Slide
%  condivise (poi non distribuite); domande finali a voce e in chat.
%  Pacchetti richiesti nel preambolo del libro:
%    \usepackage[utf8]{inputenc}, \usepackage{booktabs},
%    \usepackage{tikz} con
%    \usetikzlibrary{arrows.meta,positioning,calc,fit}
%
%  Collocazione nel libro: parte "Alcune Professioni nel Mondo
%  Videoludico", dopo AIV5_programmazione (3.7).
%  Tutte le etichette di questo capitolo hanno il prefisso "p6:".
% =====================================================================

\chapter[Il 3D artist]{Il 3D artist: specialisti, pipeline degli asset e
collaborazione}
\label{p6:cap:3d}

La presentazione di corso dei Dev Days (capitolo~\ref{p1:cap:devdays};
la trascrizione dice ``benvenuti a un'altra giornata di Database'') è quella
di \textbf{3D art}, tenuta dal coordinatore del corso, \textbf{Giacomo
Caneschi}\footnote{Nella trascrizione ``Canischi''.}. Il programma che
annuncia: un'infarinatura su chi sono i 3D artist, come li forma AIV, che cosa
fanno e che cosa \textbf{non} fanno; poi la pipeline di lavoro ``classica''
per la maggior parte degli asset 3D, per far capire anche a chi non è un
artista che cosa vuol dire fare modelli 3D da mettere in un motore ``a livello
della generazione corrente, che non è assolutamente una cosa banale''; infine
ampio spazio alle domande.

L'incontro si tiene online su Zoom (``scusate questi tre minuti di ritardo, ma
Zoom proprio non voleva collaborare: ci ho dovuto lottare un po' a mani
nude''), di mattina. La data non è detta; due indizi la collocano verso la fine
della settimana dei Dev Days, con ogni probabilità \textbf{giovedì 8 o venerdì
9 ottobre 2026}\footnote{Il relatore dice agli studenti di 3D art che li vedrà
``tra un paio di settimane'', e a uno studente che mancano ``una decina di
giorni all'inizio del corso''; le lezioni cominciano ufficialmente il 19
ottobre (capitolo~\ref{p1:cap:devdays}). È una deduzione.}.

Avvertenza di lettura: la trascrizione è automatica ma, per questo incontro,
abbastanza pulita; il relatore stesso si definisce ``logorroico'' e con ``uno
\emph{stream of consciousness} piuttosto importante''. Storpia però molti nomi
tecnici: AIV diventa ``Dive'', ``Ive'', ``Hive'', ``AIB'' e ``live'';
Allegorithmic ``l'Algoritmic''; ZBrush ``Z-Brash''; Houdini ``Udini''; high
poly ``iPoly'' o ``ai poli''; UV ``OV'' o ``V''; sculpting ``scalping''; il
lead artist ``id artist''. Dove la ricostruzione è sicura si usa la forma
corretta, altrimenti lo si dice in nota. Alcuni aneddoti e alcune affermazioni
tecniche del relatore sono inesatti o non verificabili: lo si segnala in nota.

Il percorso del capitolo:
\begin{enumerate}
  \item il relatore e la sua carriera: dal generalista allo specialista
        (sezione~\ref{p6:sec:relatore});
  \item che cos'è un 3D artist: artista e tecnico, con limiti giganteschi
        (sezione~\ref{p6:sec:cosa});
  \item conoscere il lavoro degli altri, e il lavoro a progetto
        (sezione~\ref{p6:sec:altri});
  \item il 3D artist interpreta, non inventa: il concept artist
        (sezione~\ref{p6:sec:concept});
  \item i tempi della formazione e della produzione, e perché l'arte finale
        arriva per ultima (sezione~\ref{p6:sec:tempi});
  \item le \textbf{tre macro specializzazioni}: character, environment,
        props (sezione~\ref{p6:sec:specializzazioni});
  \item la \textbf{pipeline} di un asset 3D, fase per fase
        (sezione~\ref{p6:sec:pipeline});
  \item chi comanda, il mercato, la pazienza
        (sezione~\ref{p6:sec:comanda});
  \item l'output prima di tutto, e la bellezza soggettiva
        (sezione~\ref{p6:sec:output});
  \item le domande finali (sezione~\ref{p6:sec:domande}).
\end{enumerate}

% ---------------------------------------------------------------------
\section{Il relatore: dal generalista allo specialista}
\label{p6:sec:relatore}

Caneschi è un 3D artist, ``uno sviluppatore della vecchia guardia italiana'':
sviluppa videogiochi in Italia da più di vent'anni. Nasce come \textbf{grafico
classico}: per tantissimi anni ha lavorato nell'editoria del cinema, facendo
poster, campagne pubblicitarie e cose del genere. Poi, visto che era un lavoro
che stava morendo (``lo racconterò più nello specifico ai 3D artist''), si è
spostato sul 3D, anche perché i videogiochi sono sempre stati una delle sue
passioni principali.

\subsection{Il generalista, ieri}

Lavorare nel 3D vent'anni fa era molto diverso: un lavoro ``drasticamente più
tecnico, molto meno artistico''. Si dava per scontato che chi faceva i modelli
sapesse fare anche tutto il resto: i videogiochi hanno una quarantina d'anni,
ma fino a una ventina d'anni fa farli era ``proprio magia'', e quando ci si è
resi conto che era una cosa accessibile a tecnici specializzati, si pensava che
questi tecnici fossero specializzati un po' in tutto. Per lavorare, il
relatore ha dovuto imparare non solo il 3D con perizia, ma anche ad animare, a
fare gli shader, qualche linguaggio di programmazione, e all'occorrenza fare il
designer, lavorare in engine, ``qualsiasi cosa possibile e immaginabile''.

\subsection{Lo specialista, oggi}

Per fortuna l'industria è cambiata, e si è riconosciuto (``credo che ve
l'avranno già detto in tanti'') che il \textbf{generalista}, il ``tuttologo''
che risolveva tutti i problemi, prima molto ben visto, è in realtà un
\textbf{gigantesco collo di bottiglia}. ``Più responsabilità do a una persona,
più le do la possibilità di sbagliare'', e non è detto che sappia gestirle
tutte allo stesso modo: ci sarà sempre qualcosa che le è più affine. Per questo
è arrivata alla ribalta la figura dello \textbf{specialista}: nel 3D, una via
di mezzo fra un tecnico e un artista, che deve saper fare una serie di cose, ma
che viene assunto per farne \textbf{una specifica}.

Questo è tanto più vero quanto più si sale di livello di qualità e di fondi
delle società: le grandi, i tripla A, gran parte dei ``super indie'' tendono ad
assumere specialisti. Nell'indie ``basso'', il ``garage plus'', gli schemi
saltano: ci può essere un sound designer che fa gli sprite 2D, il programmatore
che sceglie il design, ``l'artista probabilmente non si presenta nemmeno al
lavoro''. Ma se si deve avere un obiettivo, almeno all'inizio, è lavorare in
una \textbf{società di livello piuttosto alto}: lì ci si toglie soddisfazioni
vere, si trova il proprio nome in produzioni molto grandi e si cominciano a
guadagnare ``soldini veri''.

\subsection{Texture e material artist}

Più o meno nel 2015-2016, ``agli albori'' di Allegorithmic\footnote{Allegorithmic,
la società francese di Substance Painter e Substance Designer, è stata fondata
nel 2003; il relatore si riferisce con ogni probabilità all'affermazione dei
suoi software (Substance Painter esce nel 2014). È stata acquisita da Adobe nel
2019.}, ``una società che fa software per 3D artist'', il relatore ha
intrapreso il percorso da specialista: ha abbandonato il 90\% delle cose che
faceva per dedicarsi ai \textbf{materiali}. È un \textbf{texture artist},
\textbf{material artist}: il tipo di 3D artist che lavora esclusivamente ``la
parte 2D del 3D''. Un po' perché veniva dalla grafica classica e si muoveva
meglio con le texture, un po' perché gli piaceva, era un lavoro nuovo e c'era
grandissima richiesta: ha ``cavalcato l'onda''.

Gli ha portato soddisfazioni che da generalista non aveva avuto: è passato dagli
sviluppatori italiani ``zozzi'' (quelli che non fanno parte della ``rosa di 6-7
società serie'' in Italia) a lavorare con gli internazionali, con ``stipendi
giusti''; ha lavorato ``con veramente chiunque'', in Italia e all'estero. Nel
\textbf{2019} è arrivato nel tripla A, reclutato da NVIDIA e Microsoft come
uno degli sviluppatori della tecnologia RTX, il sistema di visualizzazione
basato sulla fisica; in particolare è stato uno dei texture artist ufficiali di
\textbf{Minecraft RTX}, ``il gioco che ha instaurato la tecnologia RTX nei
videogiochi correnti''\footnote{La versione RTX di \emph{Minecraft} (beta nel
2020) è stata una vetrina molto visibile del ray tracing in tempo reale, ma non
il primo gioco a usarlo: lo avevano preceduta, fra gli altri,
\emph{Battlefield~V} (2018) e \emph{Metro Exodus} (2019).}. Si definisce ``un
tizio un po' più da R\&D e un po' meno da produzione'': gli piace giocare con
cose che non è detto che funzionino. Oggi lavora anche per Amazon, sulla serie
di \emph{Fallout}, come lead (``sono un po' il capoccia''): ``vita parallela''.

% ---------------------------------------------------------------------
\section{Che cos'è un 3D artist}
\label{p6:sec:cosa}

\subsection{Un tecnico specializzato}

Il 3D artist ``non è semplicemente una persona che sta lì a pettinare le
bambole'', come si dice in gergo: è un \textbf{tecnico specializzato}, che sa
usare ``una pletora di programmi'', 12 o 13, con l'eccellenza in alcuni di
essi. Per questo, dice il relatore, la parola ``artista'' è quasi sbagliata: la
parte artistica è una parte, la tecnica è ``qualcosa di molto più
validato''\footnote{Così la trascrizione (``molto più valdato''); forse ``molto
più pesante'', ``molto più vasto''.}. Ci sono tanti artisti classici che si
danno al 3D, con un ottimo approccio alla scultura e alla pittura, ma che sono
``tecnicamente inabili'', e quindi non possono lavorare nei videogiochi.

Il videogioco è infatti \textbf{la forma di grafica 3D più complessa che
esista}. Per il cinema, la realtà aumentata, la prototipazione, la stampa 3D,
tutto è ``drasticamente molto più semplice'': gli artisti di quello che gli
sviluppatori chiamano l'\textbf{extrasettore} possono esagerare con tutto
(numero di poligoni, grandezza delle texture, qualità delle scene). Gli artisti
dei videogiochi lavorano con \textbf{limiti giganteschi}, e devono ``fare con
pochissimo, tantissimo''. Va da sé che, quando un artista di videogiochi lavora
nell'extrasettore, il suo lavoro si semplifica: per la serie di
\emph{Fallout}, dice il relatore, ``il mio lavoro è drasticamente più
semplice''.

\subsection{I limiti}

I limiti sono molteplici:
\begin{itemize}
  \item l'\textbf{hardware}, il più importante: il gioco deve girare
        idealmente a 60 fotogrammi al secondo, possibilmente di più, ed essere
        visivamente decente anche a 30. La macchina deve gestire qualcosa di
        variabile, e anche l'output degli artisti deve esserlo;
  \item l'\textbf{input del giocatore}: in ogni singolo fotogramma il giocatore
        può dare indicazioni che cambiano drasticamente tutto ciò che è a
        schermo, quindi quello che è a schermo non può essere esagerato,
        altrimenti nemmeno il computer più potente riesce a fare tutti i
        calcoli.
\end{itemize}
Bisogna quindi fare cose che sembrano fatte con tantissimi poligoni usandone
pochissimi, cose che sembrano fatte con texture 4K usando texture 2K: essere
``molto resilienti'' su come si lavora e sulle risorse che si usano.

\subsection{Artista e tecnico, da soli o in gruppo}

Il 3D artist è comunque ``assolutamente un artista'', perché produce qualcosa
di manifatturiero che esprime la sua visione e le sue capacità; ma con tanta
tecnica. Le sfumature sono molte:
\begin{itemize}
  \item ci sono 3D artist che usano più l'arte e meno la tecnica, e usano meno
        software. In generale, \textbf{più si è specializzati, meno software si
        usano}; più si è generalisti, più software servono (anche se alcune
        specializzazioni richiedono comunque molti software);
  \item ci sono artisti che lavorano \textbf{da soli} e artisti che lavorano
        \textbf{in gruppo}, e non perché uno sia ``il lone wolf'' e l'altro
        ``l'artista amico con cui vai a bere dopo'': alcune mansioni non possono
        essere affrontate in tempo utile da una persona sola, mentre in altri
        lavori è importante che si veda \textbf{una mano unica}.
\end{itemize}
I 3D artist hanno anche \textbf{capacità laterali}. Un \emph{weapon artist},
specializzazione molto richiesta all'interno dell'universo dei prop artist, sa
fare armi favolose, ma conosce anche tutto il resto: il giorno in cui sono
finite le armi da produrre, ci sarà altra roba da fare, e gli verrà assegnato
tutto ciò che gli è affine.

% ---------------------------------------------------------------------
\section{Conoscere il lavoro degli altri}
\label{p6:sec:altri}

``E questo lo dico per tutti'': un artista, ma in generale qualunque
developer, deve conoscere a grandi linee, anzi con una \textbf{profondità
media}, il lavoro di tutti gli altri sviluppatori. Un artista deve sapere che
lavoro fa un designer, che cosa gli può chiedere e che cosa no; che tipo di
programmatore ha accanto; che cosa può chiedere a un ingegnere del software, a
un tech artist, a un tech designer. Non vuol dire saper fare tutto, ``essere una
one man band'', ma avere \textbf{cultura generale sullo sviluppo dei
videogiochi}. Questa cultura si sviluppa prima di tutto \textbf{giocando}, e
poi interessandosi a tutto quello che accade ``oltre la quarta parete, oltre lo
schermo''.

\subsection{Il lavoro a progetto}

In generale, una risorsa in grado di supportare gli altri reparti, o addirittura
di sostituirli per alcune responsabilità, è molto ben vista, ed è il modo per
\textbf{tenersi il lavoro}. Perché il lavoro nei videogiochi è \textbf{a
progetto}: ``vi prendo per lavorare a \emph{Celestial Hammer}'' (un nome
inventato); il giorno in cui \emph{Celestial Hammer} è finito, ``strette di
mano, abbracci, aperitivo, si festeggia tutti insieme, e poi tutti a casa''. La
maggior parte delle persone, a casa. Restano a lavorare, e in genere fanno
carriera, quelle che si sono dimostrate \textbf{proattive}, sanno
\textbf{automotivarsi} e sanno muoversi dentro ciò di cui c'è bisogno. Quindi:
interessatevi di ciò che c'è accanto a voi, ``non pensate di essere in un vuoto:
io faccio la mia parte e via''.

\subsection{Il lavoro passa di mano}

Non è vero che un 3D artist faccia solo i modelli: bisogna pensare da dove
viene un modello e a chi va. A meno che non si faccia un gioco da soli (``se ce
la fate avete tutta la mia stima, e qualcuno ce la fa, ma in genere no''),
\textbf{il proprio lavoro passa di mano}, e su ciò che si è fatto lavora altra
gente. Sul livello creato da un level designer lavora l'environment artist, che
lo ``dressa'' e lo rende bellissimo. Quando si passa un character
all'animatore, deve essere ottimizzato per l'animazione: se ha tutti i
\emph{joint} smontati, uno scheletro brutto, o anche prima, prima del rigger,
se semplicemente non è modellato \emph{animation friendly}, ``l'animatore non
lo può animare in nessun modo'' (si veda il capitolo~\ref{p4:cap:animation}).

Il 3D artist è quindi \textbf{prima di tutto uno sviluppatore}, come tutti gli
altri ruoli; subito dopo è sia un artista sia un tecnico. E come passerà i suoi
modelli al level designer, al rigger, all'animatore, così il suo lavoro gli
viene passato da qualcuno: il \textbf{concept artist}.

% ---------------------------------------------------------------------
\section{Interpretare, non inventare: il concept artist}
\label{p6:sec:concept}

\subsection{L'arte su commissione}

Qui ci si allontana dall'idea comune di arte, che oggi fa pensare a ``un tizio
che mette un elicottero ribaltato in piazza, o una tipa che si fa tirare
addosso i pomodori''. L'arte classica nasce invece come \textbf{lavoro su
richiesta}, e così il lavoro del 3D artist. Michelangelo non faceva le sue
sculture e i suoi dipinti perché gli andava: c'era ``un prete con un sacco di
soldi'' che gli diceva ``voglio che mi fai la Pietà, voglio che mi fai il
soffitto di questo posto''\footnote{Semplificazione scherzosa: la Pietà
vaticana fu commissionata dal cardinale francese Jean de Bilhères, la volta
della Cappella Sistina da papa Giulio II.}. Allo stesso modo un 3D artist,
quando non sta facendo il portfolio, \textbf{non lavora per sé stesso}, ma per
qualcun altro; e in genere non è chiamato a inventare quasi nulla, ma a
\textbf{interpretare e ricreare nel modo più fedele possibile} le indicazioni
del concept artist.

Questo anche perché non è detto che il 3D artist venga dall'arte classica. Il
relatore stesso non è un pittore: ha ``scarabocchiato tutta la vita'', e
l'approccio più simile all'arte classica che ha avuto è stato dipingere sui
muri\footnote{Nella trascrizione ``quando dipingevo sui muri Facelbrider'':
forse ``quando facevo il writer'', cioè graffiti.}. Col 3D, ``come diciamo a
Roma, me la comando'', ma non sente di avere la libertà di inventare, perché gli
manca la cultura su tutta una serie di cose: basi anatomiche, teoria del
colore, proporzioni, stili. Per questo gli artisti 3D si appoggiano ai
\textbf{concept artist}: disegnatori formati su tutto ciò che i 3D artist non
fanno, chiamati a dare la loro visione, ``un concetto di quello che dobbiamo
andare a fare''.

\subsection{Buoni occhi, non buone mani}

Il relatore mostra due immagini, un disegno e il modello 3D ricavato: per un
occhio non abituato ai videogiochi è difficile dire quale dei due sia il
modello, e il modello non è nemmeno post-processato. È un lavoro di
\textbf{Valerio Carbone}, docente di AIV al secondo anno, ``uno dei
videoscultori più forti in Italia'': con questo modello è arrivato, se il
relatore non sbaglia, secondo a uno dei \textbf{contest di ArtStation}, ``che
non si fanno più'' e che erano ``la vetrina più importante a livello mondiale
per i 3D artist''. Da lì è iniziata la sua carriera: ora lavora su
\emph{Warhammer}\footnote{Nella trascrizione ``fa Overhammer, Tatticus'':
forse \emph{Warhammer} e un titolo legato al personaggio di Titus
(\emph{Warhammer 40,000: Space Marine}), ma non è verificabile.}, ha lavorato
``con chiunque'', ha fatto \emph{Spyro}. E la sua carriera può iniziare da
quel modello perché è \textbf{riprodotto con perizia}.

È la skill principale del 3D artist: la capacità di \textbf{leggere quello che
vede e saperlo riprodurre}. Per questo non è importantissimo aver fatto studi
d'arte classica. Un artista classico può avere un immaginario fantastico e
tantissima cultura, ma se non ha \textbf{buone mani} (per disegnare linee
morbide, per scolpire la pietra senza spaccarla) non può fare nulla. Nell'arte
digitale servono \textbf{buoni occhi}, ``e basta''; e poi memoria e
applicazione.

\subsection{Il talento non esiste}

Per questo, dice il relatore, ``è un lavoro in cui il talento non è
richiesto''. In più di vent'anni di lavoro e più di sedici di insegnamento si è
convinto che \textbf{il talento non esista}: tutto si basa sull'applicazione.
``La persona talentuosa è semplicemente una persona che ha iniziato a fare
qualcosa, ci ha messo il cuore, e con il tempo è diventata brava''. Le persone
che si avranno davanti non saranno necessariamente artisti navigati: potrebbero
aver iniziato da pochi anni, ma con una ``fissa'' talmente grande da essere
diventate molto brave.

\subsection{Studiare sempre}

Per esplorare tutto il 3D, tutte le sue applicazioni e tutti i software, non
basterebbero sette-otto anni di formazione; ma per diventare un ottimo artista e
produrre asset per la generazione corrente bastano \textbf{due-tre anni di
formazione effettiva}, e il resto si studia da soli. È importantissimo capire
che la tecnologia, soprattutto nelle immagini digitali, cambia ogni giorno:
tutto ciò che si impara oggi è vero ``per un tot di tempo''; poi potrebbe uscire
un software che rivoluziona l'industria, e lo si deve imparare. Il caso
classico: Allegorithmic ha fatto \textbf{Substance Painter} e \textbf{Substance
Designer}, di cui il relatore è specialista, che hanno cambiato radicalmente il
modo di fare le texture: prima si facevano in Photoshop, ``una fatica, un
dolore senza senso''. Quando esce un software che ``vi puzza'' di diventare lo
standard dell'industria, lo si impara subito. ``Io tuttora, dopo vent'anni,
studio''. E lo stesso vale per chiunque usi un motore o un software: il software
si aggiorna di continuo, e bisogna aggiornarsi con lui.

\subsection{Concept, key art e referenze}

Per lo \textbf{stilizzato} (lo stile ``eroico'' del modello mostrato) c'è
sempre prima il disegno di qualcuno. Per il \textbf{realismo} questa fase può
essere in parte bypassata: in genere si ha una \textbf{key art} più che una
concept art, cioè un'immagine chiave che racconta come deve essere ciò che si
farà, e poi una pletora di \textbf{referenze} per ogni singolo dettaglio. Se si
fa un character realistico, si vuole sapere il colore della barba, se ha i
brufoli, se è grosso o piccolo, di che tessuto è fatta la maglietta; se la sua
arma è nuova o rovinata, quanto, dove: nel deserto? in mezzo all'acqua? Sono
informazioni molto importanti per creare qualcosa di visivamente corretto.

% ---------------------------------------------------------------------
\section{I tempi: della formazione e della produzione}
\label{p6:sec:tempi}

\subsection{Pronti l'ultimo giorno del terzo anno}

Fra tutte le professioni che AIV forma, il 3D artist è quella che \textbf{ci
mette di più} a raggiungere il suo livello massimo: ci sono tanti software da
studiare e tanti lavori e argomenti diversi da esplorare prima che un artista
capisca davvero che cosa gli piace e che cosa gli riesce meglio. Quindi, per
tutti quelli che collaboreranno con i 3D artist: saranno perfettamente formati
\textbf{l'ultimo giorno del terzo anno}. Non si arriva al livello finale al
secondo anno per poi fare la tesi al terzo: servono tutti e tre gli anni. ``Lo
dico per i designer, per i narratori, per i programmatori'': gli artisti con cui
si comincerà a lavorare al secondo anno sono \textbf{totalmente in training} e
non hanno abbastanza conoscenze per produrre tutto ciò che serve. ``Vi dovete
dare del tempo''.

\subsection{L'arte finale arriva per ultima}

Per fortuna, nei videogiochi l'arte finale viene messa \textbf{alla fine}. Non
si crea mai una \emph{gym} né si fa debugging con i modelli finali: per un
modello come quello mostrato un professionista impiega \textbf{un mese, un
mese e mezzo}; uno studente probabilmente quattro, in cui deve fare anche tante
altre cose. Non bisogna aspettarsi la roba finita il primo giorno: nel
frattempo si sviluppa il gioco con \textbf{blockout}, modelli secondari,
\emph{dummy} (quelli base di Unreal per i character) e ``cuboni'' per
l'environment. I giochi si sviluppano così.

Tutto questo avviene in \textbf{pre-produzione}, prima di sviluppare
effettivamente il gioco; gli artisti lavorano in \textbf{produzione}: quando
il gioco è stato definito e non si fanno più cambi, si chiama l'artista per il
modello definitivo. Un cambio drastico che butta via il lavoro di un 3D artist
``è un po' una pezza'': se lo si è pagato per quattro mesi, si sono idealmente
bruciati ``10.000 euro, ve lo butto un po' così''; anche se i videogiochi hanno
spesso budget diversi, non lo si vuole mai fare.

\subsection{L'esempio di Dragon Trainer}

L'esempio che il relatore fa sempre, ``perché c'ero, sono stato coinvolto in
questa cosa'', è il film \emph{Dragon Trainer}\footnote{\emph{How to Train
Your Dragon} (DreamWorks, 2010). Che il protagonista e il drago siano stati
ripensati in corso d'opera è documentato: quando Chris Sanders e Dean DeBlois
presero in mano la regia (2008), la storia, Hiccup e Sdentato furono
profondamente rivisti, rispetto anche ai libri di Cressida Cowell, dove
Sdentato è un draghetto piccolo e comune. Non è corretto però che la decisione
sia arrivata ``a film finito'': arrivò a metà sviluppo. Il ruolo del relatore
nella produzione non è verificabile; il drago ``brutto'' della versione
originale non è identificabile con certezza dalla trascrizione (``Terrore
Verde''; nei libri c'è un drago gigantesco chiamato Green Death).}. Nella prima
versione il protagonista e il drago erano molto diversi: lui era fiero, sicuro
di sé, baldanzoso; il drago era un ``coso brutto'' ispirato alle leggende
norvegesi, con ``una serie di occhi dispari'' (``sette occhi''). Quando gli
\emph{stakeholder}, ``quelli coi soldi'', lo hanno visto, hanno detto che non
andava: le persone non avrebbero empatizzato con il protagonista, e il drago
era ``brutto come il peccato, come si dice in termine tecnico''. Sono stati
sviluppati un altro protagonista, insicuro, zoppo, con caratteristiche in cui
le persone si possono rivedere, e un drago dall'aspetto ``fico'', Sdentato,
lungo, nero, \emph{slick}.

Nei film, a fatica, si può fare: si sostituisce una serie di fotogrammi. Nei
videogiochi \textbf{non si può}: un gioco viene annunciato con enorme anticipo,
e quando è annunciato la gente ``stacca la bocca''; se non esce nell'anno
previsto, ``internet si rivolta'', ed è probabile che lo studio ``prenda delle
belle botte''. Per questo \textbf{tutte le decisioni si prendono in
pre-produzione}. In genere, quando si immagina un gioco, c'è un solo artista
che lavora con il team di design e di narrativa, il \textbf{lead artist}, che
decide una serie di cose insieme all'\textbf{art director}, il capo dei
concept artist. Loro discutono, cominciano a disegnare e a modellare, e quello
diventa la base di tutto: ``non prendo 30 artisti se non so cosa devono fare,
perché sono soldi buttati''. Gli artisti, in genere, non fanno lavoro
sperimentale: ``accendono il computer e cominciano a lavorare''.

Al tempo stesso il concept è un'immagine prospettica e basta: non c'è il retro,
non si vede che cosa c'è sotto il collo. Tutte queste parti \textbf{le deve
inventare il 3D artist}, in stile con il concept, ``perché il concept non si
muove, il modello 3D sì''.

% ---------------------------------------------------------------------
\section{Le tre macro specializzazioni}
\label{p6:sec:specializzazioni}

Gli artisti che si incontreranno in questi anni rientrano in \textbf{tre macro
categorie}, ciascuna con sottospecializzazioni: quando si incontra un artista,
idealmente vorrà fare una di queste tre cose (tabella~\ref{p6:tab:macro}).

\begin{table}[htbp]
\centering
\small
\begin{tabular}{@{}lp{3.3cm}p{3.3cm}p{3.3cm}@{}}
\toprule
 & \textbf{Character} & \textbf{Environment} & \textbf{Props} \\
\midrule
Chi lo fa & pochi artisti per studio; una sola mano (salvo clothing e
grooming) & molti artisti, sempre in team, con level designer, programmatori,
tech artist & tutti, e in particolare i junior: è il primo lavoro \\
\addlinespace
Tempi & un character ogni 2-3 mesi (stilizzato: circa un mese per un
professionista) & 5-6 mesi per un environment; un modulo in giornata &
tantissimi asset: 200 in due mesi e mezzo, circa 4 al giorno \\
\addlinespace
Risorse & circa 300.000 poligoni a schermo in un tripla A & un muro può essere
due triangoli & spesso comprati in pack, soprattutto nel realismo \\
\addlinespace
Carattere & ``un po' più di ego'' & bisogna ``uccidere il proprio ego'' &
responsabilità minori, ma un prop bello è difficile \\
\bottomrule
\end{tabular}
\caption{Le tre macro specializzazioni del 3D artist come le descrive il
relatore.}
\label{p6:tab:macro}
\end{table}

\subsection{I character}

I personaggi sono \textbf{il cuore} di quasi ogni videogioco (tranne quelli
basati sulle interfacce; ``anche i giochi di carte hanno i personaggi, ormai''),
perché serve un personaggio per identificarsi e immergersi. A volte non si vede
(si è in prima persona, magari è solo raccontato, si vedono solo le mani): ma
allora quelle mani devono comunicare tantissimo del personaggio. I character
sono la cosa più \emph{glamour} che può fare un 3D artist: il \textbf{90\%}
degli studenti che entrano in AIV dice ``voglio fare i personaggi''; solo il
\textbf{20\%} circa di quelli che escono lo dice ancora.

Perché fare personaggi è un \textbf{lavoro molto lungo}. Un semplice personaggio
stilizzato può richiedere un mese a un professionista, per farlo bene, dopo i
feedback; a un artista in formazione, probabilmente il doppio. Un character
realistico non è più complesso, ma molto più lungo: nella stilizzazione i
dettagli vengono rimossi, alcuni esasperati (``pompati''), e solo quelli
principali vanno riportati; nel realismo bisogna metterci tutto. Anzi: ``il
realismo è palloso e non lo facciamo mai''; si fa l'\textbf{iperrealismo}, il
realismo pompato, definendo il personaggio fino ai \textbf{pori della pelle}:
si texturizza per far capire se in una zona i pori sono grassi o no, se c'è un
po' di peluria. Quindi, quando serve un character, bisogna essere consci che
l'artista ci metterà parecchio tempo, soprattutto se lo si vuole fatto bene.

Anche perché quell'artista deve fare tutto, o il 90\% delle cose, \textbf{da
solo}: i character devono sembrare fatti dalla stessa mano, con dettagli
congrui, senza distacchi di qualità o di stile, altrimenti si finisce nella
\textbf{uncanny valley}, dove non si capisce che cosa sia stilizzato e che cosa
realistico, ``né carne né pesce''. E gli artisti hanno qualità diverse: in due
sulla stessa cosa, uno sarà più bravo, e ci saranno parti più belle di altre.
Nei personaggi tutto deve sembrare estremamente consistente. Le due sole cose
che di solito si appaltano ad altri artisti sono:
\begin{itemize}
  \item il \textbf{clothing}, solo nei personaggi realistici: vestiti fatti con
        la fisica, da artisti specializzati;
  \item il \textbf{grooming}: peli, capelli, barbe, pellicce, piume. È un lavoro
        a parte, che si fa in modo molto diverso dalla modellazione classica;
        se si ha un character artist bravo nel grooming ``si è trovato un
        tesoro'', altrimenti si chiama un \emph{groomer}, ``una persona che di
        lavoro pettina le bambole''. È identico a come si fa nel cinema: nel
        cinema i peli e le piume sono molti di più e più definiti, ma il
        workflow è lo stesso.
\end{itemize}
Il character artist è anche ``quello con un po' più di ego'', che ci mette
parecchio del suo: in genere è la persona ``un po' più spinosa'' con cui avere a
che fare. Qui ci si avvicina all'arte classica, e si sente il prodotto molto
proprio.

\paragraph{Low poly e high poly.} Il modello che si vede nel gioco non è mai
fatto come lo si vede. Quello che si muove nel motore è una forma geometrica
rozza, il \textbf{low poly}; la versione bellissima, l'\textbf{high poly},
non può girare in un motore. ``Anche se quelli di Epic vi dicono che, grazie a
Nanite, potete buttare i vostri modelli high poly dentro Unreal e farli girare,
è assolutamente falso''\footnote{È un'affermazione netta: Nanite permette
effettivamente di usare in Unreal Engine 5 mesh statiche con un numero di
poligoni molto superiore a quello tradizionale, ma non elimina la necessità di
ottimizzare, e per anni non ha supportato le mesh animate con scheletro
(\emph{skeletal mesh}), supporto introdotto solo di recente in forma
sperimentale. Il relatore ci torna più avanti
(sezione~\ref{p6:sec:pipeline}).}: un high poly ``vi cuoce il computer'',
perché può essere fatto di 70 milioni di poligoni, mentre un character a
schermo in un tripla A è fatto di circa \textbf{300.000 poligoni}, ``e sono già
tantissimi''. Per arrivare ai pori della pelle bisogna però aggiungere dettagli
su dettagli: come si fa, se a schermo non ci si possono mettere milioni di
poligoni? Si fanno in modo \textbf{simulato}, con un software apposito che
permette di simulare anche miliardi di poligoni e poi di trasformarli in
poligoni veri. Quindi il character artist sbozza il modello, fa i dettagli, fa
il low poly, fa le texture: tutto il processo da solo.

\subsection{Gli environment}

Lavoro drasticamente diverso è quello dell'\textbf{environment artist}, che
\textbf{riempie i livelli dei level designer} con una serie di modelli. È un
lavoro che si fa \textbf{sempre in team}: i livelli sono estremamente grandi,
anche quelli piccoli sono grandi scene 3D, e il lavoro di una persona non basta
mai. Gli environment artist lavorano da soli solo quando fanno il portfolio;
altrimenti sempre in gruppo, con i level designer, e molto anche con i
programmatori e con i tech artist. È un lavoro molto più corale, in cui si può
gestire tutto l'aspetto del 3D o solo alcune parti: c'è chi fa solo luci
(\emph{light design}), chi solo effetti speciali e post-produzione, chi solo
moduli architettonici (città, castelli); o si può saper fare tutto ma essere
presi per una cosa sola.

Per fare environment bisogna ``un attimo \textbf{uccidere il proprio ego}'',
perché il proprio lavoro deriva da quello dei level designer: bisogna
rispettarne la visione e il lavoro. È un lavoro più aperto al feedback, anche
perché la visione da rispettare è in genere quella del lead artist, che
``bene o male saprà come funziona''. Ed è un mondo \textbf{più tecnico}: se un
character fa 300.000 poligoni, ``un muro, spesso e volentieri, dentro un
videogioco sono due triangoli''. Si lavora con ancora meno risorse, si deve
creare roba leggera che sembri estremamente bella, e saper lavorare
l'ottimizzazione (``le sfumature sono infinite, ma non è oggi il momento di
parlarne''). In compenso i moduli si fanno molto più velocemente: ``una sedia
la fai in giornata'', anche in un'oretta, ad alcuni servono due giorni. Sembra
un lavoro meno importante e meno complesso, ma proprio perché coinvolge un
sacco di persone è estremamente complesso. In uno studio servono \textbf{pochi
character artist} e \textbf{molti environment artist}.

\subsection{I props}

Il terzo ruolo, che può avere un artista dedicato ma che in genere si chiede di
saper fare a tutti, ed è il primo lavoro che si fa quando si comincia a fare 3D
per lavoro, sono i \textbf{props}: l'oggettistica, ``tutta''. Barili, formaggi,
ma anche le pistole e le spade sono props; alcune parti di un'armatura possono
esserlo, e quindi non le fa il character artist; un portafoglio, un cellulare,
una catena, le scarpe. È il lavoro che si dà ai \textbf{junior}. In realtà
fare un prop bello è complesso: bisogna \textbf{raccontare la storia di
quell'oggetto}, e ``forse è più difficile raccontare la storia di una
bottiglietta d'acqua che la storia di una persona''. Semplicemente gli oggetti
sono più piccoli e più semplici: c'è differenza fra un'idra a nove teste e ``un
barile, che è fondamentalmente un cilindro ciccione''. Le responsabilità sono
minori, anche perché tantissimi props oggi \textbf{si comprano}: un grande pack
di props anni Ottanta con dentro il mangianastri, la cassetta, il pupazzo di
He-Man\footnote{Nella trascrizione ``il pupazzo dei men'': con ogni
probabilità He-Man, o i Masters of the Universe.}, già pronti.

Va ricordato anche nei propri progetti: nello \textbf{stilizzato}, avendo scelto
uno stile specifico, in genere tutto ciò che si vede a schermo è stato fatto da
qualcuno interno allo studio; nel \textbf{realismo}, spesso le cose vengono
dall'\emph{outsourcing} o da pack. AIV ha avuto per molto tempo una
collaborazione con \textbf{Quixel}, quella di Megascans, ``il primo software
che utilizzava scansioni del mondo vero per farci le texture''\footnote{Più
precisamente Megascans è una libreria di scansioni fotogrammetriche (oggi di
Epic Games), affiancata da software di Quixel come Mixer e Bridge; non è stata
l'unica né forse la prima, ma è stata la più diffusa.}: il relatore aveva una
libreria di scansioni e materiali, che poi ha ritrovato dentro \emph{Call of
Duty}, ``non dentro il gioco dell'indie zozzo''. È legale usare materiale
scansionato e asset comprati. D'altra parte per un artista è importante fare
lavoro proprietario, e negli studi più importanti ultimamente si cerca di fare
tutto internamente, per avere visual più unici: ``se tutti usano gli stessi
props, poi i giochi sono tutti uguali''.

Questo è ancora più vero perché si usano sempre più \textbf{sistemi
procedurali} e sempre meno lavoro fatto a mano (lo stilizzato è ancora in gran
parte a mano). Nel procedurale non si dipinge il bronzo a mano: si prende un
materiale bronzo e lo si associa a un modello, e il materiale ha dei valori per
l'ossidazione, la superficie, i bozzi, le impurità. Il difetto è che, se non lo
si lavora abbastanza, \textbf{è tutto uguale}, perché tutti hanno accesso alle
stesse risorse: se tutti fanno la stessa spada, rischiano di sembrare fatte
tutte dello stesso materiale.

La grande differenza fra i props e le altre specializzazioni è la
\textbf{quantità}. A un character artist si può chiedere un character ogni due
o tre mesi; al team di environment, un environment in cinque-sei mesi. Un prop
artist deve produrre tanto: quando il relatore ne ha assunto uno per un
progetto di cui era lead, gli ha fatto una proposta che prevedeva \textbf{più
di 200 asset in due mesi e mezzo}, cioè circa quattro modelli finiti al giorno,
cinque giorni alla settimana (``il povero cristiano poi ce l'ha fatta''). Fra un
sasso e il boss finale di \emph{Dark Souls}\footnote{Nella trascrizione ``il
boss finale di Endsols''.} la differenza di complessità è grande, anche se
vanno fatti tutti per bene: chi fa cose più semplici ne fa di più, ma le
responsabilità sono le stesse.

\subsection{Fare soldi, e far uscire i giochi}

Il videogioco è \textbf{un bene condiviso}, a cui tutti lavorano insieme perché
venga bello: si lavora per fare soldi, ma anche perché quei soldi siano
associati a una \textbf{credibilità lavorativa}. Un amico del relatore lavora da
tanti anni per Warner Bros.\ e guadagna benissimo, ma gli hanno chiuso quattro
progetti ed è disperato, perché non è riuscito a far uscire nulla: ``ha farmato
un sacco di soldi, ma al pubblico non c'è niente: è come se non avesse fatto
nulla''.

Le specializzazioni sono molte, e non serve conoscerle tutte oggi: se ne
parlerà al corso, e la scelta si fa da soli alla fine del percorso. Gli artisti
penseranno di fare una delle tre cose, ma possono fluttuare: si comincia
volendo fare character, poi si fa environment, si scopre che è bellissimo e si
cambia idea.

% ---------------------------------------------------------------------
\section{La pipeline di un asset 3D}
\label{p6:sec:pipeline}

``Che fa questo povero 3D artist tutta la vita? Impazzisce davanti a una serie
di software''. Il relatore commenta la pipeline lineare di produzione di un
asset 3D (figura~\ref{p6:fig:pipeline}).

\begin{figure}[htbp]
\centering
\begin{tikzpicture}[>=Stealth, font=\scriptsize, node distance=0.35cm,
  box/.style={draw, rounded corners, align=center, text width=2.1cm,
              minimum height=0.9cm}]
  \node[box, fill=gray!15] (c) {concept\\(pre-produzione)};
  \node[box, right=of c] (l) {low poly /\\blockout};
  \node[box, right=of l] (s) {high poly\\(inorganico) o\\sculpting (organico)};
  \node[box, right=of s] (r) {retopology\\(solo dopo lo\\sculpting)};
  \node[box, below=0.9cm of r] (u) {UV layout};
  \node[box, left=of u] (b) {baking\\(normal map\\e altre mappe)};
  \node[box, left=of b] (t) {texturing};
  \node[box, left=of t, text width=2.4cm] (o) {inorganico $\to$ level design\\organico $\to$ rig $\to$ animazione};
  \draw[->] (c) -- (l); \draw[->] (l) -- (s); \draw[->] (s) -- (r);
  \draw[->] (r) -- (u); \draw[->] (u) -- (b); \draw[->] (b) -- (t); \draw[->] (t) -- (o);
  \node[below=0.35cm of $(o.south)!0.5!(u.south)$, font=\scriptsize\itshape]
    {nel frattempo i designer sviluppano il gioco con il blockout (low poly)};
\end{tikzpicture}
\caption{La pipeline lineare di un asset 3D. Il UV layout si fa dopo la
retopology se c'è sculpting, altrimenti prima del dettaglio; il baking non
compariva sulla slide e il relatore lo aggiunge a voce. Ricostruzione
schematica.}
\label{p6:fig:pipeline}
\end{figure}

\subsection{La fase concettuale}

C'è sempre una \textbf{fase concettuale}, che viene in genere \textbf{dopo}
che i narrative e gli scrittori hanno scritto la base della storia, che sono
state create le prime \emph{gym} e il documento di design\footnote{Nella
trascrizione ``il God Book di design'': con ogni probabilità il GDD
(\emph{game design document}) o la ``bibbia'' del progetto.}. Definite le cose
scritte, si traducono in immagini: si fa il concept del gioco, con immagini che
rappresentano la visione superiore del gioco, e si comincia a definire come
modellare i modelli. Il \emph{level concept} va in parallelo alle prime fasi di
sviluppo del 3D, perché prima viene tutto ciò che riguarda il gioco in sé, poi
i visual: un livello non è solo visual, è prima di tutto \textbf{giocabilità},
spazi, metriche. ``Poi dopo lo facciamo diventare bello''.

\paragraph{Il level designer non fa l'enrichment.} Attenzione, dice il
relatore: vi potrà essere detto che ci sono level designer che fanno il
cosiddetto \emph{enrichment}. ``Questa è una roba che non è assolutamente vera
in tantissime società'', soprattutto nella parte alta. Il level designer si
occupa esclusivamente di come funziona il livello e dei suoi ingombri, non della
forma finale (si veda il capitolo~\ref{p2:cap:gamedesign}). ``Dire io voglio
quella pianta proprio lì non è un lavoro da level designer'': il level designer
ci mette un cilindro e dice ``da lì non si passa'', oppure ``è un cilindro sopra
un piano, da lì non si passa, ma c'è una cosa che si muove col vento''. Punto.
Il \textbf{designer} deve fare una cosa che \textbf{funziona}, l'\textbf{artista}
una cosa \textbf{bella}, e il gioco deve fondere le due cose. Non si possono
sviluppare le due skill contemporaneamente in tempo utile: prima il level
design, poi l'arte. In realtà si lavora a braccetto, si decide un po' insieme;
ma ogni figura ha la sua skill, e per i visual ci si affida alla mano
dell'artista.

I concept si fanno in pre-produzione; qualcosa di meno importante finisce in
produzione, ma le cose principali (il personaggio principale, l'aspetto
generale del mondo, i mostri principali) si fanno tutte prima di produrre
effettivamente il gioco.

\subsection{Low poly e blockout: si sviluppa con i cuboni}

Poi l'artista fa il \textbf{low poly}: un modello estremamente rozzo che
comprende la forma generale e soprattutto gli \textbf{ingombri}, importantissimi
(``se ti serve un gigante e io ti do un nano, è un problema''): un gigante con
la sua scala, fatto per bene, con pochissimi poligoni, super leggero. Sarà la
base che si dettaglierà fino alla forma finale; ed è la forma ``cubosa, non
definita'' con cui \textbf{lavorano i designer}.

Il relatore lo dice perché gestisce anche lo sviluppo dei videogiochi in AIV,
l'iniziativa \textbf{Game Dev}\footnote{Dal contesto, il progetto
multidisciplinare dal secondo anno in poi (i Dev Projects,
sezione~\ref{p1:sec:devprojects}).}: \textbf{non servono i modelli finali per
sviluppare il gioco}; lo si fa tutto con i cuboni. I modelli finali servono
prima delle \emph{milestone} importanti, delle consegne. Se tutto gira con il
modello rozzo, basta sostituirlo: se è statico, è immediato; se si muove,
bisogna fare il \emph{retargeting} delle animazioni, ``una roba rognosa, ma non
infattibile, anzi, lo fanno tutti''. Con quei modelli bianchi, senza texture, i
designer e chi lavora in engine testano tutto: post-produzione, illuminazione.

\subsection{UV layout}

Intanto i 3D artist, dopo il low poly, devono ``tagliare'' il modello per
prepararlo alla fase successiva: si ottiene l'\textbf{UV layout}, per poter
dipingere in due dimensioni quello che sarà il materiale, la texture, l'aspetto
del modello. Il 3D non è dipingibile direttamente, ``perché il 3D sono dati e le
texture sono pixel'': l'UV layout lega la posizione dei pixel in 2D alle
coordinate del modello 3D. ``È una roba pallosa che va fatta per forza. Più la
fate bene, più vi vorranno bene i vostri lead''. Se c'è la retopology, l'UV
layout si fa dopo; se non c'è, prima del dettaglio. Il modello viene
``spezzettato'' e steso su un piano, trasformato in un'immagine; ciò che si
dipinge sull'immagine finisce sul modello.

\subsection{Il detailing: high poly o sculpting}

Poi c'è il \textbf{detailing}, in due modi:
\begin{itemize}
  \item per tutto ciò che è \textbf{inorganico} (una macchina, una pistola) si
        fa l'\textbf{high poly modeling}: si prende il low poly e si aggiungono
        dettagli, dettagli, dettagli, fino alla forma finale;
  \item per tutto ciò che è \textbf{organico} (in genere un personaggio, una
        creatura) si fa lo \textbf{sculpting}: non si modella in modo
        geometrico, ma si prende ``una creta virtuale di poligoni simulati'' e
        con la penna si scava, si disegna, si aggiunge e si toglie massa, fino
        alla forma finale.
\end{itemize}
Il software di sculpting mostrato è \textbf{ZBrush}, ``un software che vi fa
uscire il sangue dagli occhi quando guardate l'interfaccia'': ``un miliardo di
menu'', quando oggi i software usano drag and drop, tasto destro, icone; tanto
che spesso la prima lezione è come personalizzarne l'interfaccia per renderla
più comprensibile. Lo sculpting è una fase lunga, che gli artisti tendono ad
allungare all'infinito perché è la più bella, ``la più sugosa'', quella in cui
si vede l'aspetto effettivo del modello. Ma va \textbf{chiusa al più presto}: un
professionista la fa fra una settimana e dieci giorni al massimo, perché entro
un mese deve consegnare; uno studente può passarci tranquillamente un mese, un
mese e mezzo.

\subsection{La retopology}

La \textbf{retopology} c'è solo quando c'è lo sculpting. Per un modello
inorganico si passa dal low poly all'high poly, ed è finita. Nella scultura,
mentre si lavora il dettaglio, gli ingombri e le forme possono cambiare
leggermente: quindi c'è una fase in più, in cui \textbf{si rifà il low poly
sopra l'high poly}. All'inizio una cosa estremamente rozza, poi scolpita, poi
ricreata low poly, ``ma non proprio rozza'', per farla girare nel motore.

Sulla slide la topologia di un volto è colorata a zone: rappresentano i
\textbf{gruppi muscolari} che si muovono insieme, ed è tutto lavoro fatto per
l'animazione. Si potrebbero mettere i poligoni come si vuole e ottenere più o
meno lo stesso aspetto; ma proprio perché gli artisti, soprattutto i character
artist, non lavorano mai da soli, bisogna lavorare in un certo modo: altrimenti,
quando si assegnano le \emph{blend shape} (la tecnica con cui si animano le
facce), la faccia non reagisce come dovrebbe (si veda la discussione sulla
topologia nel capitolo~\ref{p4:cap:animation}).

\subsection{Il baking}

Il \textbf{baking} non è sulla slide, ma il relatore lo aggiunge: è il
trasferimento di tutti i dettagli dell'high poly o della scultura sul low poly,
in modo da simularli nel motore. Il giocatore vede l'high poly, ma nel gioco
gira il low poly\footnote{La trascrizione dice ``girano 20.000 poligoni ma io
vedo 20.000'': evidentemente un numero è trascritto male (il senso è ``ne
vedo milioni'').}. I dettagli vengono trasportati in una texture da assegnare al
low poly: la mappa più importante è la \textbf{normal map}, ``la mappa di
asperità'', che rappresenta i piccoli cambiamenti di altezza della
superficie\footnote{La trascrizione dice ``i cambiamenti di altezza da uno a 5
cm'': forse un ordine di grandezza indicativo dei dettagli che conviene
affidare alla normal map invece che alla geometria.}, che così non serve
modellare; insieme alle altre mappe permette un livello di dettaglio
altissimo con pochissime risorse. ``Se in un \emph{Total War} ci devono essere
500 modelli a scena, devono essere leggeri''.

Esistono tecnologie che rendono tutto più leggero, come \textbf{Nanite} di
Unreal; ma a Nanite, dice il relatore, ``non possono essere assegnate le
skeletal mesh, quindi niente che si muove'': ``è una tecnologia un po' ibrida,
che ancora non sappiamo bene che faccia''\footnote{Si veda la nota sulla
sezione dei character: il supporto di Nanite alle mesh animate è stato
introdotto di recente in forma sperimentale.}.

\subsection{Il texturing}

Infine il \textbf{texturing}, la specializzazione del relatore: lo standard è
\textbf{Substance Painter}, che AIV insegna (sulla slide, ``con l'interfaccia
drasticamente modificata''). È il dare un aspetto materico ai modelli, ``la
loro pelle, il loro colore, le specifiche dei materiali'':
\begin{itemize}
  \item nello \textbf{stilizzato} si parla di volumi, valori di colore,
        gradienti;
  \item nel \textbf{realismo} si parla delle controparti fisiche: non solo il
        colore, che ``in verità non esiste: esiste il colore percepito'', ``una
        roba che si sono inventati gli artisti per riprodurre il mondo''; ma
        anche la microsuperficie, la metallicità, le zone di
        auto-ombreggiamento, per generare una texture realistica.
\end{itemize}
Finito il texturing, \textbf{il modello è finito}. Se è inorganico, passa in
genere ai level designer; se è organico, al team di animazione, dove prima si
sviluppa il rig e poi le animazioni.

\paragraph{Un 3D artist anima? No.} ``No, assolutamente no. È un altro lavoro''.
Si possono conoscere pillole di animazione, saper fare un po' di rigging, ma
non è richiesto (``io ho dovuto impararlo, ai tempi, e ancora soffro'').

\subsection{I software}

Una slide mostra i software usati per ogni fase, ``e non sono neanche tutti'':
per il blockout, Maya, 3ds Max o Blender (se ne usa uno dei tre); lo stesso per
la retopology; per il texturing può capitare di dover usare tutti i software
insieme per un risultato migliore. Il corso non ha nemmeno il tempo di insegnare
tutto di ogni software, anche perché non serve: molte cose di Maya servono per
render iperrealistici, per animazioni, per il CAD con le spline, e non vengono
insegnate. Ma restano ``tòtere'' di software complessi. Il relatore lo dice
soprattutto ai narrative e ai game designer che impareranno un software come
Unreal: ``i grafici sono passati nelle vostre scarpe, e l'hanno dovuto fare
dieci volte: invece di imparare un software, ne hanno dovuti imparare
dodici''.

% ---------------------------------------------------------------------
\section{Chi comanda, il mercato, la pazienza}
\label{p6:sec:comanda}

\subsection{Nessuna figura ha la precedenza}

Il lavoro del 3D artist passa per molte mani, quindi bisogna essere
\textbf{molto flessibili} e aperti alla collaborazione con chi viene prima e
con chi viene dopo. Non si va dal concept artist a dirgli che una cosa è brutta:
``se piace al lead artist, la vostra opinione non conta''. Non si va a ``rompere
le scatole'' ai level designer, che devono fare il loro lavoro: ci si lavora
insieme, in collaborazione totale. \textbf{Non c'è una figura che ha la
precedenza su un'altra}; l'unica che comanda sono gli \textbf{stakeholder},
``quelli che ci mettono i soldi: arrivo io, metto 10 milioni, comando io''.

E l'idea può venire da chiunque. Il relatore racconta sempre che il concept di
\emph{Horizon Zero Dawn} è stato deciso da un tech artist, ``cioè da un
programmatore'', che ha portato un prototipo di dinosauro robot fatto con i
cuboni, che veniva colpito e perdeva i pezzi\footnote{L'aneddoto non è
verificabile nella forma riportata; è noto che il concept del gioco nacque da
una proposta interna a Guerrilla Games, scelta fra diverse candidate.}. ``Non
c'era la visione incredibile di un narratore, non c'erano le skill infinite di
un game designer''. Le figure che AIV forma \textbf{sviluppano} l'idea, non
devono per forza averla. Chi comanda? \textbf{I tre lead insieme}: il lead
designer, il lead artist e il lead developer, che comandano chi sta sotto di
loro; e loro fanno quello che dicono l'art director, il technical director e,
prima ancora, gli stakeholder (si veda anche la
sezione~\ref{p2:sec:gerarchia}). ``In questo settore non fai mai quello che ti
pare''.

\subsection{I giochi si fanno per gli altri}

A meno di non fondare una società e diventare capi di sé stessi; ``allora lì
puoi fare quello che ti pare, e probabilmente la tua società fallirà e tu
finirai nelle peste''. Perché anche da capi di sé stessi bisogna studiare il
\textbf{mercato} e le \textbf{nicchie}: i videogiochi non si fanno mai per noi,
ma per chi ci gioca. Fare quest'anno un roguelike di carte sarebbe ``da
folli'': è uscito \emph{Slay the Spire 2}, e per chi gioca ad altri deck builder
ne sono usciti di bellissimi (il relatore ne cita uno\footnote{Nella
trascrizione ``Shroom Gloom'': il titolo non è verificabile.}). Il relatore è
appassionato di carte (gioca a \emph{Magic}, colleziona Pokémon), ma non fa un
gioco di carte ``perché mi va'': c'è il mercato, ci sono prima di tutto i
giocatori. Anche un modello 3D si fa perché piaccia al pubblico, non per sé:
``altrimenti vieni licenziato''.

L'aneddoto: un concept artist, per un gioco ``targettizzato 12-14''\footnote{Così
nella trascrizione (``12x14''): forse una fascia d'età 12-14 anni, o una
classificazione PEGI 12.}, stilizzato, ``pupazzoso'', con violenza non
realistica, aveva inserito nei concept delle armi una serie di riferimenti
sessuali. Il relatore gli ha rimandato i concept con tutte quelle parti
cerchiate, spiegando il target; il concept artist ha risposto ``no, sono io che
mi immagino questo mondo, e questo mondo è fatto così''. ``Un abbraccio,
grazie: da oggi non lo decidi più tu''. Strette di mano, contratto stracciato.

\subsection{Pazienza, ma con dei limiti}

``Il clou, la cosa più importante che ho da dire oggi a tutti i non 3D artist'':
nel collaborare con i 3D artist, \textbf{soprattutto finché non vi pagano},
bisogna avere \textbf{pazienza}. Ci vuole tempo, ci vuole un po' di
\emph{trial and error}. ``Il tizio ancora non me l'ha consegnato'' è normale;
poi, a un certo punto, ci sono dei limiti oltre i quali non è più normale. Nel
lavoro, in quel caso, si licenzia: un videogioco, se va bene, venduto a 9 euro
rischia di fare milioni; se ci sono persone che lavorano male o sono lente,
quei milioni rischiano di non arrivare. Per questo le si riprende una volta, due
volte, e alla terza le si allontana: o si viene licenziati brutalmente, o si
viene ``dimensionati''\footnote{Il relatore usa il termine ``white quitting'';
il fenomeno descritto si chiama di solito \emph{quiet firing}.}: si viene
assegnati a cose sempre più stupide e ignobili, finché si dice ``ma che ci
resto a fare qua dentro? me ne vado''.

% ---------------------------------------------------------------------
\section{L'output prima di tutto}
\label{p6:sec:output}

Il lavoro del 3D artist è, in fondo, semplice: ``parto da una roba rozza e la
faccio sempre più bella''. Ma ha infinite sfumature e workflow differenti. Ci
vuole quindi flessibilità quando si chiede qualcosa a un 3D artist: non è detto
che la sappia fare, né che la faccia con un workflow classico. ``L'importante,
per tutti, è l'\textbf{output}''.

È una cosa che il relatore ha imparato in Ubisoft: ha due amici fraterni nelle
risorse umane e una volta l'anno va a Parigi a parlare con gli artisti che,
come lui, sono stati formati da Allegorithmic per insegnare quei software (``poi
Allegorithmic è stata comprata da Adobe, tristemente le cose sono un po'
cambiate''). Gli artisti di Ubisoft dicono che ai loro art director interessa
una sola cosa: \textbf{che quello che fai sia bello}. Come l'hai fatto non
importa, purché l'abbia fatto tu e non l'abbia rubato a qualcuno. Il processo è
molto personale, perché esistono tanti software per fare la stessa cosa: per il
texturing lo standard è Painter, ma si può usare 3D-Coat per lo stilizzato o
Mari per i volti iperrealistici. Quindi non ci si chiede mai come il proprio
artista stia facendo una cosa: conta che abbia un buon ritmo, che consegni
quando deve e che ciò che consegna sia bello.

\subsection{La bellezza è soggettiva}

E la bellezza è soggettiva. Ci sono giochi con una grafica, per il relatore,
``di una bruttezza incredibile'' che hanno avuto un successo enorme. Cita sempre
\emph{Kenshi}, ``esteticamente di una bruttezza catastrofica'', con una grafica
da primi anni Duemila: associata a un gioco sistemico con un design favoloso,
``fa il giro'', diventa adatta, ``quasi bella''. In questi anni c'è il ritorno
dell'8 bit, dei \emph{boomer shooter} (giochi che richiamano \emph{Doom}): a
volte una grafica vecchio stile è quasi richiesta. La prima volta che ha visto un
modello di \emph{Valheim} ha pensato ``ma che sta schifezza'': rozzo,
texturizzato con un'unica mappa per tutto il gioco, a pixel, ``una roba di un
team che non aveva artisti''. Poi ci ha giocato: non solo è un gioco favoloso, ma
a livello artistico ``tutto torna'', quei pochi pixel ``fanno scopa con gli altri
pixel che gli stanno accanto''. Oggi la \textbf{prodezza tecnica} non si valuta
quasi più, ``la valuti su GTA, sui giochi che paghi 70 euro''; altrimenti c'è
l'accettazione di ciò che si ha, se piace. ``Le cose brutte fanno il giro e
diventano belle'': un qualcosa va sempre giudicato \textbf{nel suo insieme},
perché anche una serie di modelli brutti può fare un gioco bello.

\subsection{Le slide non si danno}

Il relatore non distribuisce le slide: sono una versione ridotta di quelle che
darà agli artisti il primo giorno, e a chi non è un artista servono fino a un
certo punto; le usa solo come guida visiva. ``L'importante è che mi avete
ascoltato''. Ai 3D artist: ``ci vediamo tra un paio di settimane''; agli altri:
``dovrete aspettare un annetto e mezzo, ma tanto ci vediamo tutti a Game Dev''.

% ---------------------------------------------------------------------
\section{Le domande finali}
\label{p6:sec:domande}

\subsection{Per un gioco 2D ci si affida a un 3D artist?}

Mentre sta salutando arriva una domanda, ``molto importante, di cui forse avrei
dovuto parlare'': Federico chiede se, per un gioco 2D, ci si affidi comunque a
un 3D artist. ``No, assolutamente no'': è un'altra skill, serve un
\textbf{grafico 2D}. Anche l'animazione 2D e quella 3D sono lavori
drasticamente diversi, non solo per una dimensione in meno: nel 2D interessano
movimenti fluidi lungo una linea, ``su una fascia''; nel 3D ogni vertice del
modello è gestito da tre curve, e spesso per armonizzare il movimento bisogna
ritoccarle, decidere se in un punto la curva fa un angolo acuto o una
sinusoide.

L'animazione, aggiunge, è mezza tecnica e mezza artistica: la parte tecnica è il
\textbf{rigging}, quella artistica sono le \textbf{capacità attoriali}, perché
l'animatore deve far recitare il proprio modello (capitolo~\ref{p4:cap:animation}).
In un gioco 2D conta solo la recitazione; in un gioco 3D la tecnica è tutto. Un
gioco 2D lo potrebbe fare anche un concept artist: in genere c'è un
illustratore, e a seconda della tecnica si usano il \emph{ragdoll} o le clip di
animazione. AIV non forma 2D artist: fa ``la parte 2D del 3D'', il texturing.

E le \textbf{interfacce}? Le fa un \textbf{UI artist}\footnote{Nella
trascrizione ``un Wii Artist''.}, che in genere è un 2D artist; nei giochi
futuri degli studenti potrebbe farle un 3D artist, o un designer, ``chiunque
abbia un po' di gusto''. In genere c'è un grafico classico, abile a fare
``pulsantino spento, pulsantino acceso'' in modo gustoso. In una jam o in un
progettino chiunque può imparare ad animare in 2D (``non è difficilissimo;
farlo bello è molto, molto, molto difficile''); ma i giochi che si faranno in
AIV è bene farli \textbf{in 3D}, perché AIV forma 3D artist, e chi lavorerà sui
visual farà 3D: ``se vi fate fare le cose in 2D è un problema''.

\subsection{Quando si comincia a usare la tavoletta grafica?}

Marco, presumibilmente un futuro artista, chiede quando si inizia a usare la
tavoletta grafica e quale consigliare. La seconda domanda va rifatta alla prima
lezione (al relatore, se si è a Roma o online; ``a Lorenzo'' se si è
nell'\emph{Experience} di Milano), o in privato su Discord, perché ``è troppo
lunga spiegarla a tutti''. Si comincia a usarla da subito. Non tutti gli artisti
la usano: il relatore, che fa materiali procedurali, non la usa mai; ma la si
insegna subito, perché qualcuno potrebbe avere affinità con la scultura, che si
fa \textbf{solo con la tavoletta}. Fra le centinaia di studenti avuti, uno
``psycho'' scolpiva col mouse ed era pure ``mezzo bravo'', ma era un
\emph{unicum}. Le domande specifiche sul 3D, da artista ad artista, vanno fatte a
lezione, dove ci sarà tempo, e dove il relatore racconterà di nuovo la sua vita,
più nel dettaglio.

\subsection{Quando comincia a lavorare il concept artist?}

In \textbf{pre-produzione}: è il tipo di artista che lavora prima dei 3D artist.
In pre-produzione la prima cosa che si vuole a livello visivo è la visione
generale, e i primi concept li fa in genere l'\textbf{art director} (il
relatore dice prima ``il lead artist'' e si corregge), che durante la produzione
non disegna mai, ma in pre-produzione deve anche disegnare un po'. Più si va
verso i ruoli di leadership, meno si fa ``manovalanza'', e per un art director
disegnare un concept è manovalanza: scriverà una serie di cose e farà degli
sketch.

I concept più importanti si fanno in pre-produzione: il personaggio principale,
le \emph{mood board} delle cose più importanti dei livelli, armi, oggetti,
nemici. Perché quando si parte in produzione bisogna avere materiale su cui far
lavorare il team di produzione, cioè grafici, animatori e programmatori. Poi si
va ``a cascata'': il concept dell'ultimo livello si può fare in produzione,
perché verrà fatto alla fine; quello del primo livello va fatto in
pre-produzione, perché chi arriva a fare il livello deve avere il concept
pronto, altrimenti è rallentato. È un punto chiave: ogni volta che questo ciclo
non viene rispettato ci sono grossi problemi (``ah, ma non c'è il concept''), e
bisogna aspettare, magari mentre i designer hanno trovato una feature nuova e
devono testarla subito. La parte visual si testa anche tutti insieme, ognuno
dice la sua, poi decidono gli artisti e l'art director: è una cosa condivisa, in
cui si chiede l'impressione delle persone al proprio livello. ``Il lead artist
non abbraccia l'ultimo dei grafici per chiedergli se gli piace'':
probabilmente lo presenterà al lead developer e al lead designer: ``ho preso la
tua idea e ci ho fatto questo. Secondo te funziona? Se lo facciamo in tempo
utile, diventiamo tutti poveri e moriamo prima, oppure riusciamo a chiuderlo,
farci dei soldi e prosperare?''.

\paragraph{Rispettare la pipeline.} Se la pipeline non viene rispettata, in
genere tocca \textbf{tornare indietro}, ``e tornare indietro sulla produzione è
veramente una roba brutta''. Fare un videogioco è complesso: più ci si avvicina
alla consegna, più tutto dovrebbe diventare leggero, ``la gente non ha le
crisi''; tornare indietro è molto negativo. Si fa spessissimo, ma bisogna cercare
di non farlo. Per questo la pre-produzione è così importante, anche per il
singolo artista e per qualsiasi tecnico: la \textbf{fase preparatoria} è chiave.
``Non dovete mai partire sull'onda dell'entusiasmo'': quando arriva un lavoro, la
prima cosa è \textbf{documentarsi}. Nell'arte servono i concept o una bella
ricerca; ma vale anche per un designer che volesse fare un gioco ambientato in
un ospedale senza sapere che cosa succede in un ospedale: ne uscirebbe un gioco
``palloso, terribile'', che non rispecchia la realtà.

\paragraph{Il pitch e la tech demo.} Si può anche prototipare senza concept
(il relatore si perde nel discorso, ``si vede che sono uno che lavora'', e poi
lo ritrova). Quando si presenta un \textbf{pitch}, spesso lo presentano solo i
game designer, parlando con tutti, ma in sostanza da soli: si va a cercare i
soldi proprio per prendere grafici, programmatori, animatori, quindi non si hanno
concept né disegni. ``Adesso si fanno con l'IA; prima si scaricava un po' di roba
online'' e si mettevano nelle presentazioni le immagini più vicine all'idea.
Insieme al pitch spesso c'è una \textbf{tech demo}, valida solo per il pitch:
magari verrà buttata quando si farà il gioco, oppure diventerà la chiave per
svilupparlo. I concept vengono dopo il pitch. L'ideazione viene sempre prima e la
può fare chiunque, ma poi serve impostare una pipeline corretta, ``altrimenti
non si consegna''. Tutti i giochi che il relatore ha visto fallire avevano gestito
male la pipeline di produzione: si erano arenati, avevano grandi colli di
bottiglia, mancava documentazione. ``Con un'ottima documentazione fai un'ottima
fase preparatoria, e in genere non puoi fallire'': magari si arriva un po' più
lunghi, con qualche imprevisto, ma raramente si fallisce. E vale per tutto ciò
che riguarda un prodotto, anche il marketing.

\subsection{Un muro che si distrugge: che cosa fa il programmatore?}

Un programmatore chiede: per creare un muro che si distrugge in più parti, deve
fare qualcosa con i modelli, o si applicano proprietà che simulano la fisica in
automatico? Ha visto video su Blender in cui bastava applicare qualcosa al
modello. ``Se c'è una roba che dici: la metto e succede qualcosa, c'è un
programmatore che l'ha scriptata: il codice c'è, semplicemente tu non l'hai
visto''. Le vie sono due:
\begin{itemize}
  \item il muro è \textbf{preparato per rompersi}, cioè già fratturato: con un
        algoritmo di \emph{fracturing} che lo rompe in forme realistiche, o
        rotto a mano; poi si implementa la fisica per farlo crollare;
  \item il muro si rompe \textbf{dinamicamente}: e allora ``tutto quello che
        hai visto su Blender è copiato da \textbf{Houdini}'', dove non serve
        codice ma bisogna essere un artista super tecnico, un \textbf{tech
        artist}, ``una persona che non pettina le bambole, ma a cui piace vedere
        le cose che funzionano''. Houdini è un'altra categoria, un altro lavoro
        rispetto alla modellazione.
\end{itemize}
Il relatore sconsiglia di far rompere le cose nei primi progetti: gli algoritmi
di fracturing sono ottimo materiale da portfolio per programmatori, proprio
perché sono complessi. ``Se il muro non si deve rompere, non lo rompere''; se
proprio si deve, servono tutti, il grafico e il programmatore. ``L'altro giorno
mi dicevano: sarebbe bellissimo un gioco in cui con la spada ti taglio il
braccio, e nel punto in cui lo taglio dinamicamente si stacca. Fermo, fermo: non
è fattibile''. Girerebbe ``su computer da NASA, quelli con cui farmavano i
bitcoin e ora addestrano le intelligenze artificiali'', non ``sul computer del
disperato come noi''\footnote{L'affermazione è troppo netta: il taglio dinamico
in un punto qualsiasi è stato realizzato in tempo reale, per esempio nella
``Blade Mode'' di \emph{Metal Gear Rising: Revengeance} (2013), anche se con
molte semplificazioni; ma è certamente un sistema complesso, fuori portata per
un progetto studentesco.}.

\subsection{Che cosa posso fare prima che inizi il corso?}

Marco chiede se, mancando una decina di giorni all'inizio, ci sia qualcosa da
fare per prepararsi, o se sia meglio aspettare per non prendere brutte
abitudini. Per prendere brutte abitudini bisogna fare una cosa molte volte, e
dieci giorni non bastano; ma sono arrivate in AIV persone autodidatte che
facevano tutto nel modo sbagliato, quindi il consiglio è \textbf{aspettare}.
``Divertiti, fai tutto quello che ti va, perché poi tocca fare i seri'': il
corso di 3D art è molto intenso, perché ci sono tanti software da insegnare.
``Rilassati, respira, fai altre cose e preparati a imparare insieme a noi''.

\subsection{I saluti}

Il relatore deve chiudere l'aula (``serve per una riunione''). Le domande da
artista ad artista si possono fargli in privato, anche prima dell'inizio delle
lezioni; tutti gli altri, per domande strette sul 3D, possono scrivergli su
Discord, ``sto su Discord appositamente''. Dall'anno prossimo, con i progetti di
Game Dev, ci si vedrà spesso. ``Per la terza volta vi ringrazio, vi saluto, vi
auguro buona giornata, e ci vediamo presto''.
